\chapter{Processi aleatori}

Un segnale che arriva a un ricevitore non è un singolo numero, e nemmeno un vettore di pochi numeri: è una successione
di valori che continua nel tempo. Il rumore in un circuito, la sequenza di bit trasmessa da un modem, la temperatura
misurata ogni giorno, il prezzo di un'azione ogni minuto sono tutte \textbf{successioni aleatorie}: non si possono
prevedere esattamente, ma hanno delle regolarità statistiche. Lo strumento per descriverle è il \textbf{processo
aleatorio}, un esperimento il cui risultato non è un numero ma un'intera successione. In questo capitolo, l'ultimo del
secondo pacchetto di slide, ne diamo la definizione nel caso del tempo discreto, vediamo come si caratterizza e
introduciamo la nozione più importante, la \textbf{stazionarietà}: l'idea che le proprietà statistiche del processo non
cambino nel tempo.

% ============================================================
\section{Richiami: la caratterizzazione completa}
% ============================================================

Si consideri uno spazio di probabilità arbitrario $(\Omega, S, P)$: $\Omega$ è lo spazio campione, $S$ è una
$\sigma$-algebra di eventi di $\Omega$ (Capitolo~4) e $P : S \to [0, 1]$ è una legge di probabilità. Ricordiamo che una
variabile aleatoria è una funzione
\[ X : \omega \in \Omega \;\longmapsto\; X(\omega) \in \mathcal{X} \subseteq \R , \]
discreta se tale è l'alfabeto $\mathcal{X}$, continua se tale è $\mathcal{X}$. Una coppia di variabili aleatorie è
l'applicazione $(X, Y) : \omega \in \Omega \longmapsto (X(\omega), Y(\omega)) \in \mathcal{X} \times \mathcal{Y}
\subseteq \R^2$, e una $m$-pla, ovvia generalizzazione della coppia, si scrive come vettore aleatorio (Capitolo~16):
\[ \vett{X}(\omega) = [X_1(\omega), X_2(\omega), \ldots, X_m(\omega)]^T \in \mathcal{X}_1 \times \mathcal{X}_2 \times
   \cdots \times \mathcal{X}_m \subseteq \R^m . \]
Una variabile, una coppia o un vettore si dicono \term{completamente caratterizzati} (statisticamente) quando se ne
conosce la legge, in una delle forme equivalenti della tabella: la CDF, che esiste sempre, oppure la pmf nel caso
discreto o la PDF nel caso continuo. La PDF si ottiene derivando la CDF una volta rispetto a ciascuna variabile.

\begin{center}
\renewcommand{\arraystretch}{1.9}
\begin{tabular}{@{}lccc@{}}
  \toprule
  & CDF (sempre) & pmf (caso discreto) & PDF (caso continuo) \\ \midrule
  $X$ & $F_X(x) = P(X \leq x)$ & $p_X(x) = P(X = x)$ & $f_X(x) = \dfrac{\mathrm{d}}{\mathrm{d}x} F_X(x)$ \\
  $(X, Y)$ & $F_{X,Y}(x, y) = P(X \leq x, Y \leq y)$ & $p_{X,Y}(x, y) = P(X = x, Y = y)$ &
  $f_{X,Y}(x, y) = \dfrac{\partial^2 F_{X,Y}(x, y)}{\partial x\, \partial y}$ \\
  $\vett{X}$ & $F_{\vett{X}}(\vett{x}) = P(X_1 \leq x_1, \ldots, X_m \leq x_m)$ & $p_{\vett{X}}(\vett{x}) = P(X_1 = x_1,
  \ldots, X_m = x_m)$ & $f_{\vett{X}}(\vett{x}) = \dfrac{\partial^m F_{\vett{X}}(\vett{x})}{\partial x_1 \cdots
  \partial x_m}$ \\ \bottomrule
\end{tabular}
\end{center}

% ============================================================
\section{La regola della catena per i vettori aleatori}
% ============================================================

Nel Capitolo~8 abbiamo visto la regola della catena per tre variabili discrete. Per i processi ci servirà con un numero
qualsiasi di variabili.

\teorema[catenavettori]{Regola della catena per i vettori aleatori}{
  Per un vettore aleatorio discreto $\vett{X} = [X_1, \ldots, X_m]^T$:
  \[ \begin{aligned}
     p_{\vett{X}}(\vett{x}) = {} & P(X_1 = x_1)\, P(X_2 = x_2 \mid X_1 = x_1)\, P(X_3 = x_3 \mid X_1 = x_1, X_2 = x_2) \\
     & \times \cdots \times P(X_m = x_m \mid X_1 = x_1, \ldots, X_{m-1} = x_{m-1}) ,
     \end{aligned} \]
  cioè, in forma compatta,
  \[ p_{\vett{X}}(\vett{x}) = \prod_{i=1}^{m} p_{X_i \mid X_1, \ldots, X_{i-1}}(x_i \mid x_1, \ldots, x_{i-1}) , \]
  dove per $i = 1$ si conviene che $p_{X_1 \mid X_0}(x_1 \mid x_0) = p_{X_1}(x_1)$ (non c'è nessuna condizione). Per i
  vettori continui vale lo stesso con le densità:
  $f_{\vett{X}}(\vett{x}) = \prod_{i=1}^{m} f_{X_i \mid X_1, \ldots, X_{i-1}}(x_i \mid x_1, \ldots, x_{i-1})$, con
  $f_{X_1 \mid X_0}(x_1 \mid x_0) = f_{X_1}(x_1)$.
}

\begin{dimostrazione}[si stacca una variabile alla volta]
\passo{La probabilità composta} Si parte da $P(A \cap B) = P(A \mid B)\, P(B)$, la legge della probabilità composta del
Capitolo~4 (che le slide chiamano qui legge di Bayes). Con $A = \{X_2 = x_2, \ldots, X_m = x_m\}$ e
$B = \{X_1 = x_1\}$:
\[ p_{\vett{X}}(\vett{x}) = P(X_2 = x_2, \ldots, X_m = x_m \mid X_1 = x_1)\, P(X_1 = x_1) . \]
\passo{Si ripete dentro il condizionamento} La probabilità condizionata a $\{X_1 = x_1\}$ è una probabilità
(Teorema~\ref{teo:condprob}), quindi si può staccare allo stesso modo la variabile $X_2$:
\[ p_{\vett{X}}(\vett{x}) = P(X_3 = x_3, \ldots, X_m = x_m \mid X_1 = x_1, X_2 = x_2)\, P(X_2 = x_2 \mid X_1 = x_1)\,
   P(X_1 = x_1) . \]
\passo{Si itera} Ogni volta una variabile passa dal gruppo «da calcolare» al gruppo «condizionante»; dopo $m - 1$ passi
il gruppo da calcolare contiene solo $X_m$, e si ottiene il prodotto dell'enunciato.
\passo{Il caso continuo} Si ripete lo stesso ragionamento con la definizione di densità condizionata del Capitolo~14,
«congiunta uguale condizionata per marginale», applicata ogni volta a un gruppo di variabili.
\end{dimostrazione}

La regola della catena descrive il vettore \textbf{un pezzo alla volta}, nell'ordine in cui le variabili «escono»:
prima $X_1$, poi $X_2$ sapendo $X_1$, poi $X_3$ sapendo le prime due, e così via. È il modo naturale di descrivere un
segnale che si svolge nel tempo, campione dopo campione. Due casi particolari li conosciamo già.
\begin{itemize}
  \item Se le variabili sono \textbf{indipendenti}, ogni condizionata coincide con la marginale e la regola diventa il
        prodotto delle marginali, $p_{\vett{X}}(\vett{x}) = \prod_i p_{X_i}(x_i)$.
  \item Nella gaussiana bivariata del Capitolo~16 abbiamo trovato $f_{X_1 X_2}(x_1, x_2) = f_{X_1}(x_1)\,
        f_{X_2 \mid X_1}(x_2 \mid x_1)$, con la marginale $\mathcal{N}(0, 1)$ e la condizionata
        $\mathcal{N}(\rho x_1, 1 - \rho^2)$: è la regola della catena con $m = 2$.
\end{itemize}

\begin{esempio}{Tre assi di fila}
Da un mazzo di 52 carte se ne estraggono tre, senza reinserimento; sia $X_i = 1$ se la $i$-esima carta è un asso e
$X_i = 0$ altrimenti. Qual è la probabilità di estrarre tre assi? Le variabili non sono indipendenti (ogni asso estratto
rende meno probabile il successivo), ma con la regola della catena ogni fattore è facile:
\[ P(X_1 = 1, X_2 = 1, X_3 = 1) = \underbrace{\frac{4}{52}}_{P(X_1 = 1)} \cdot \underbrace{\frac{3}{51}}_{P(X_2 = 1
   \mid X_1 = 1)} \cdot \underbrace{\frac{2}{50}}_{P(X_3 = 1 \mid X_1 = 1, X_2 = 1)} = \frac{24}{132\,600} \approx
   1{,}8 \cdot 10^{-4} . \]
Ogni condizionata si calcola guardando il mazzo rimasto: dopo un asso restano 3 assi su 51 carte, dopo due assi 2 su 50.
\end{esempio}

% ============================================================
\section{Processi aleatori a tempo discreto}
% ============================================================

\dfn{Processo aleatorio a tempo discreto}{
  Si definisce \term{processo aleatorio tempo-discreto} un'applicazione che a ogni elemento dello spazio campione fa
  corrispondere una successione:
  \[ X : \omega \in \Omega \;\longmapsto\; \{X(n, \omega)\}_{n \in \Z} , \]
  dove $\Z$ è l'insieme dei numeri interi relativi (positivi, negativi e zero), che fanno da istanti di tempo.
}

Una variabile aleatoria associa a ogni risultato dell'esperimento un numero; un processo aleatorio gli associa
un'\textbf{intera successione} di numeri, uno per ogni istante. È quindi una funzione di due variabili, $n$ e $\omega$,
e la si può guardare in due modi.
\begin{itemize}
  \item \textbf{Fissato il risultato $\omega$} (in orizzontale nella figura): $X(n, \omega)$ è una successione di numeri
        ben precisi, non più aleatoria, che si chiama \term{realizzazione} del processo. È il segnale che si osserva
        davvero in una prova, per esempio quello che compare sullo schermo di un oscilloscopio. Per ogni $\omega$ il
        processo si realizza in una sequenza di valori che stanno in un certo intervallo (anche infinito).
  \item \textbf{Fissato l'istante $n = n_0$} (in verticale): al variare di $\omega$, $X(n_0, \omega)$ assume valori
        diversi, uno per ogni elemento dello spazio campione: è una \textbf{variabile aleatoria}. È come «campionare
        verticalmente» il processo: si guarda il valore all'istante $n_0$ in tutte le realizzazioni possibili.
\end{itemize}
Come per le variabili aleatorie, d'ora in poi si omette $\omega$ e si scrive semplicemente $X(n)$.

Per fissare le idee useremo spesso la \textbf{sinusoide con fase a caso}, $X(n) = A\cos(2\pi f_0 n + \Phi)$, dove la
fase $\Phi$ è scelta a caso, uniformemente su un periodo. Ogni realizzazione è una sinusoide perfetta: il caso decide
soltanto da dove parte. È il modello di un oscillatore che si accende in un istante qualsiasi.

\begin{figure}[H]
\centering
\begin{tikzpicture}
\begin{groupplot}[group style={group size=1 by 4, vertical sep=0.3cm, xticklabels at=edge bottom}, realizzazione,
  width=15.5cm, height=3cm, xmin=-0.8, xmax=23.8, ymin=-1.3, ymax=1.3, ytick={-1,1}, xtick={0,5,10,15,20}]
\nextgroupplot
  \fill[evid] (axis cs:4.55,-1.3) rectangle (axis cs:5.45,1.3);
  \node[anchor=south, font=\small, text=mspiB!80!black] at (axis cs:5,1.3) {$X(n_0)$, con $n_0 = 5$};
  \addplot[mspiA] coordinates {(0,-0.99) (1,-0.79) (2,-0.38) (3,0.14) (4,0.61) (5,0.93) (6,0.99) (7,0.79) (8,0.38) (9,-0.14) (10,-0.61) (11,-0.93) (12,-0.99) (13,-0.79) (14,-0.38) (15,0.14) (16,0.61) (17,0.93) (18,0.99) (19,0.79) (20,0.38) (21,-0.14) (22,-0.61) (23,-0.93)};
  \addplot[mspiB, mark size=2.2pt] coordinates {(5,0.93)};
  \node[anchor=west, font=\small] at (axis cs:24.3,0) {$X(n, \omega_1)$};
\nextgroupplot
  \fill[evid] (axis cs:4.55,-1.3) rectangle (axis cs:5.45,1.3);
  \addplot[mspiA] coordinates {(0,0.35) (1,0.77) (2,0.99) (3,0.94) (4,0.64) (5,0.17) (6,-0.35) (7,-0.77) (8,-0.99) (9,-0.94) (10,-0.64) (11,-0.17) (12,0.35) (13,0.77) (14,0.99) (15,0.94) (16,0.64) (17,0.17) (18,-0.35) (19,-0.77) (20,-0.99) (21,-0.94) (22,-0.64) (23,-0.17)};
  \addplot[mspiB, mark size=2.2pt] coordinates {(5,0.17)};
  \node[anchor=west, font=\small] at (axis cs:24.3,0) {$X(n, \omega_2)$};
\nextgroupplot
  \fill[evid] (axis cs:4.55,-1.3) rectangle (axis cs:5.45,1.3);
  \addplot[mspiA] coordinates {(0,0.97) (1,0.96) (2,0.70) (3,0.25) (4,-0.27) (5,-0.72) (6,-0.97) (7,-0.96) (8,-0.70) (9,-0.25) (10,0.27) (11,0.72) (12,0.97) (13,0.96) (14,0.70) (15,0.25) (16,-0.27) (17,-0.72) (18,-0.97) (19,-0.96) (20,-0.70) (21,-0.25) (22,0.27) (23,0.72)};
  \addplot[mspiB, mark size=2.2pt] coordinates {(5,-0.72)};
  \node[anchor=west, font=\small] at (axis cs:24.3,0) {$X(n, \omega_3)$};
\nextgroupplot[xlabel=$n$]
  \fill[evid] (axis cs:4.55,-1.3) rectangle (axis cs:5.45,1.3);
  \addplot[mspiA] coordinates {(0,-0.25) (1,-0.70) (2,-0.96) (3,-0.97) (4,-0.72) (5,-0.27) (6,0.25) (7,0.70) (8,0.96) (9,0.97) (10,0.72) (11,0.27) (12,-0.25) (13,-0.70) (14,-0.96) (15,-0.97) (16,-0.72) (17,-0.27) (18,0.25) (19,0.70) (20,0.96) (21,0.97) (22,0.72) (23,0.27)};
  \addplot[mspiB, mark size=2.2pt] coordinates {(5,-0.27)};
  \node[anchor=west, font=\small] at (axis cs:24.3,0) {$X(n, \omega_4)$};
\end{groupplot}
\end{tikzpicture}
\caption{Quattro realizzazioni della sinusoide con fase a caso, con $A = 1$ e $f_0 = \frac{1}{12}$ (un periodo ogni 12
campioni). In orizzontale, ogni riga è una realizzazione: una successione di numeri. In verticale, i valori all'istante
$n_0 = 5$ (in rosso, nella fascia gialla) sono quattro determinazioni della variabile aleatoria $X(5)$.}
\end{figure}

La variabile aleatoria $X(n_0)$ ha una sua densità, che in generale dipende dall'istante $n_0$. La si indica con
$f_{X(n)}(x; n)$, dove il punto e virgola separa la variabile $x$ dall'istante $n$, che fa da parametro.

\dfn{Stazionarietà del primo ordine}{
  Un processo aleatorio si dice \term{caratterizzato al primo ordine} se ne è nota la PDF $f_{X(n)}(x; n)$ per ogni $n$.
  Se questa PDF non dipende dall'istante considerato, $f_{X(n)}(x; n) = f_X(x)$ per ogni $n$, il processo si dice
  \term{stazionario del primo ordine}: in tal caso basta assegnare un'unica PDF del primo ordine.
}

\begin{esempio}{La sinusoide con fase a caso è stazionaria del primo ordine}
Scriviamo la fase come $\Phi = 2\pi U$, con $U \sim \mathcal{U}\left( -\frac12, \frac12 \right)$: così $\Phi$ è
uniforme su un periodo intero. Fissato $n$, $X(n) = A\cos(2\pi U + \varphi_n)$ con $\varphi_n = 2\pi f_0 n$: è
esattamente la variabile della sinusoide con fase a caso del Capitolo~12, con la fase $\varphi = \varphi_n$. Là avevamo
trovato la densità dell'arcoseno, $\frac{1}{\pi\sqrt{A^2 - x^2}}$ per $\abs{x} < A$, che \textbf{non dipende da
$\varphi$}. Quindi
\[ f_{X(n)}(x; n) = \frac{1}{\pi\sqrt{A^2 - x^2}} \qquad \text{per ogni } n : \]
il processo è stazionario del primo ordine. In parole: se non si sa da dove parte la sinusoide, nessun istante è diverso
dagli altri.
\end{esempio}

Il primo ordine descrive ogni istante da solo, ma non dice nulla dei legami fra istanti diversi: se il segnale in un
istante è alto, sarà alto anche all'istante dopo? Per rispondere serve la descrizione congiunta delle coppie di
campioni.

\dfn{Stazionarietà del secondo ordine}{
  Un processo aleatorio si dice \term{caratterizzato al secondo ordine} se è nota la PDF congiunta
  $f_{X(n_1), X(n_2)}(x_1, x_2; n_1, n_2)$ per ogni coppia di istanti $n_1$, $n_2$. Si dice \term{stazionario al secondo
  ordine} se, per qualsiasi intero $h$,
  \[ f_{X(n_1), X(n_2)}(x_1, x_2; n_1, n_2) = f_{X(n_1 + h), X(n_2 + h)}(x_1, x_2; n_1 + h, n_2 + h) . \]
}

In altre parole, in un processo stazionario al secondo ordine la caratterizzazione congiunta di due campioni dipende
solo dalla loro \textbf{distanza temporale}, non dalla loro posizione: la PDF congiunta non cambia se i due istanti
vengono spostati «rigidamente», insieme, in anticipo o in ritardo.
\begin{center}
\begin{tikzpicture}[x=0.66cm]
  \draw[-{Stealth}, mspigray] (-0.6,0) -- (20.6,0) node[right, black, font=\small] {$n$};
  \foreach \i in {0,...,20} \draw[mspigray] (\i,-0.08) -- (\i,0.08);
  \fill[mspiA] (3,0) circle (2.6pt) node[below=3pt, font=\scriptsize, black] {$n_1$};
  \fill[mspiA] (7,0) circle (2.6pt) node[below=3pt, font=\scriptsize, black] {$n_2$};
  \fill[mspiB] (12,0) circle (2.6pt) node[below=3pt, font=\scriptsize, black] {$n_1 + h$};
  \fill[mspiB] (16,0) circle (2.6pt) node[below=3pt, font=\scriptsize, black] {$n_2 + h$};
  \draw[{Stealth}-{Stealth}, mspiA, line width=0.7pt] (3,0.4) -- node[above, font=\scriptsize] {distanza $n_2 - n_1$} (7,0.4);
  \draw[{Stealth}-{Stealth}, mspiB, line width=0.7pt] (12,0.4) -- node[above, font=\scriptsize] {stessa distanza} (16,0.4);
  \draw[-{Stealth}, mspigray, dashed, line width=0.7pt] (5,1.0) to[bend left=18]
    node[above, font=\scriptsize, black] {traslazione rigida di $h$: la PDF congiunta non cambia} (14,1.0);
\end{tikzpicture}
\end{center}
Un processo stazionario al secondo ordine lo è anche al primo: la densità di un solo campione si ottiene integrando la
congiunta rispetto all'altro (Capitolo~14), e se la congiunta non cambia con la traslazione non cambia nemmeno il suo
integrale. Il viceversa, invece, non è vero: due processi con le stesse densità del primo ordine possono legare in modo
molto diverso i campioni vicini.

% ============================================================
\section{Caratterizzazione completa e stazionarietà in senso stretto}
% ============================================================

Il primo e il secondo ordine guardano uno o due istanti alla volta. La descrizione completa li guarda tutti insieme,
in qualsiasi numero.

\dfn{Caratterizzazione completa di un processo}{
  Un processo aleatorio $X(n)$ si dice \term{completamente caratterizzato} (statisticamente) se, detto $M$ un intero
  arbitrario e detti $n_1, n_2, \ldots, n_M$ $M$ istanti di tempo arbitrari, il vettore aleatorio
  \[ \vett{X} = [X(n_1), X(n_2), \ldots, X(n_M)]^T \]
  ha una densità di probabilità $f_{\vett{X}}(x_1, x_2, \ldots, x_M; n_1, n_2, \ldots, n_M)$ nota.
}

\dfn{Stazionarietà in senso stretto}{
  Un processo aleatorio si dice \term{stazionario in senso stretto di ordine $M$} se la sua densità di ordine $M$ è
  invariante per traslazione:
  \[ f_{\vett{X}}(x_1, \ldots, x_M; n_1, \ldots, n_M) = f_{\vett{X}}(x_1, \ldots, x_M; n_1 + h, \ldots, n_M + h)
     \qquad \forall\, h . \]
  Si dice \term{stazionario in senso stretto} se lo è per ogni ordine $M$.
}

\teorema[ordini]{La stazionarietà passa agli ordini più bassi}{
  Un processo stazionario di ordine $M$ lo è anche di qualunque ordine $i \leq M$.
}

\begin{dimostrazione}
\passo{} Si fissano $i$ istanti $n_1, \ldots, n_i$ e se ne aggiungono altri $M - i$ qualsiasi, $n_{i+1}, \ldots, n_M$. La
densità di ordine $i$ si ottiene da quella di ordine $M$ integrando rispetto alle $M - i$ variabili in più
(marginalizzazione). Spostando tutti gli istanti di $h$, la funzione da integrare non cambia (stazionarietà di ordine
$M$), e quindi non cambia nemmeno il risultato dell'integrale, la densità di ordine $i$.
\end{dimostrazione}

\dfn{Processo indipendente}{
  Un processo si dice \term{indipendente} (o \term{a campioni indipendenti}) se, comunque si scelgano un intero $M$ e
  gli istanti $n_1, \ldots, n_M$, il vettore $\vett{X} = [X(n_1), \ldots, X(n_M)]^T$ è costituito da variabili
  aleatorie indipendenti, cioè
  \[ f_{\vett{X}}(x_1, \ldots, x_M; n_1, \ldots, n_M) = f_{X(n_1), \ldots, X(n_M)}(x_1, \ldots, x_M; n_1, \ldots, n_M)
     = \prod_{i=1}^{M} f_{X(n_i)}(x_i; n_i) . \]
}

Per un processo indipendente, quindi, la caratterizzazione completa si riduce a quella del primo ordine: basta
conoscere la densità di ogni singolo campione.

\subsection{Processi a valori discreti}

Si chiama \term{processo ad ampiezza discreta} (o, per brevità, \term{processo discreto}) un processo le cui realizzazioni
sono sequenze di numeri che possono assumere valori in un alfabeto discreto. Tutte le definizioni viste per i processi a
valori continui si estendono a questi processi, con la sola differenza che al posto delle densità si usano le pmf. Per
esempio, un processo discreto è stazionario in senso stretto se, comunque si scelgano $M$, gli istanti
$n_1, \ldots, n_M$ e uno spostamento $h$,
\[ P(X(n_1) = x_1, \ldots, X(n_M) = x_M) = P(X(n_1 + h) = x_1, \ldots, X(n_M + h) = x_M) . \]
Ovviamente anche qui la stazionarietà di ordine $M$ implica quella di ogni ordine $i \leq M$; in particolare un processo
stazionario ha una pmf marginale indipendente dal tempo.

\begin{esempio}{Il processo di Bernoulli e il processo quaternario}
Un caso di notevole importanza è il \term{processo indipendente binario}, con $X(n) \in \{-1, 1\}$ e
$P(X(n) = -1) = P(X(n) = 1) = \frac12$, detto anche \term{processo di Bernoulli} (in generale le probabilità possono
essere $p$ e $1 - p$): è il modello di una sequenza di bit trasmessi con i livelli $\pm 1$, come quella del Capitolo~15.
Un altro esempio è il processo \term{quaternario} indipendente, con $X(n) \in \{-2, -1, 1, 2\}$ e livelli
equiprobabili, $P(X(n) = i) = \frac14$ per $i \in \{-2, -1, 1, 2\}$: ogni simbolo trasporta due bit.
\begin{center}
\begin{tikzpicture}
\begin{groupplot}[group style={group size=1 by 2, vertical sep=0.5cm, xticklabels at=edge bottom}, realizzazione,
  width=15.5cm, height=3.3cm, xmin=-0.8, xmax=29.8, xtick={0,5,10,15,20,25}]
\nextgroupplot[ymin=-1.4, ymax=1.4, ytick={-1,1}]
  \addplot[mspiA] coordinates {(0,-1) (1,-1) (2,-1) (3,1) (4,1) (5,1) (6,1) (7,1) (8,-1) (9,-1) (10,-1) (11,-1) (12,-1) (13,-1) (14,1) (15,-1) (16,-1) (17,1) (18,1) (19,-1) (20,1) (21,-1) (22,1) (23,1) (24,-1) (25,1) (26,1) (27,-1) (28,1) (29,-1)};
  \node[anchor=west, font=\small] at (axis cs:30.3,0) {Bernoulli};
\nextgroupplot[ymin=-2.5, ymax=2.5, ytick={-2,-1,1,2}, xlabel=$n$]
  \addplot[mspiC] coordinates {(0,2) (1,1) (2,-2) (3,2) (4,-1) (5,-2) (6,-1) (7,-1) (8,1) (9,-2) (10,1) (11,-2) (12,1) (13,-1) (14,-1) (15,1) (16,1) (17,1) (18,1) (19,-1) (20,1) (21,-2) (22,1) (23,1) (24,-2) (25,-2) (26,1) (27,-2) (28,2) (29,1)};
  \node[anchor=west, font=\small] at (axis cs:30.3,0) {quaternario};
\end{groupplot}
\end{tikzpicture}
\end{center}
Una realizzazione di ciascuno dei due processi (simulazione).
\end{esempio}

\teorema[indipstaz]{Processi discreti indipendenti e stazionari}{
  Se un processo discreto è stazionario e indipendente, per ogni scelta di $M$, degli istanti $n_1, \ldots, n_M$ e dei
  valori $x_1, \ldots, x_M$:
  \[ P(X(n_1) = x_1, X(n_2) = x_2, \ldots, X(n_M) = x_M) = \prod_{i=1}^{M} P(X(n_i) = x_i) = \prod_{i=1}^{M}
     p_X(x_i) , \]
  dove $p_X(x)$ è la pmf marginale, che non dipende dal tempo. Viceversa, un processo indipendente che ha la stessa pmf
  marginale in ogni istante è stazionario in senso stretto.
}

\begin{dimostrazione}
\passo{La formula} La prima uguaglianza è l'indipendenza; la seconda la stazionarietà, per cui la marginale di ogni
campione è la stessa pmf $p_X$.
\passo{Il viceversa} Se il processo è indipendente e ha sempre la stessa marginale, vale la stessa formula, e il prodotto
$\prod_i p_X(x_i)$ non contiene gli istanti $n_i$: spostandoli tutti di $h$ non cambia nulla.
\end{dimostrazione}

Per un processo discreto indipendente e stazionario, dunque, l'intera caratterizzazione completa si riassume in
\textbf{una sola pmf}. Per esempio, nel processo di Bernoulli la probabilità di osservare la sequenza $+1, +1, -1$ in tre
istanti qualsiasi è $\left( \frac12 \right)^3 = \frac18$, e nel processo quaternario la probabilità che due campioni
valgano entrambi 2 è $\frac14 \cdot \frac14 = \frac{1}{16}$, dovunque si trovino.

% ============================================================
\section{Processi stazionari in senso lato}
% ============================================================

Conoscere la caratterizzazione completa di un processo è quasi sempre troppo: servirebbero le densità congiunte di un
numero qualsiasi di campioni. Come per le variabili e i vettori, ci si accontenta spesso di una caratterizzazione
sintetica, fatta di medie e correlazioni. Ci concentriamo sui processi a tempo discreto, ma quanto diremo vale anche per
quelli a tempo continuo.

\teorema[medie2]{Medie di un processo stazionario del secondo ordine}{
  Se $X(n)$ è stazionario del secondo ordine:
  \begin{enumerate}[label=\alph*)]
    \item la media $E[X(n)] = \mu_X(n) = \mu_X$ non dipende da $n$;
    \item la media del prodotto $E[X(n_1) X(n_2)]$ dipende solo da $\abs{n_2 - n_1}$, e non separatamente da $n_1$ e
          $n_2$.
  \end{enumerate}
}

\begin{dimostrazione}[nel caso discreto; nel continuo si usano le densità]
\passo{a)} Per la stazionarietà $P(X(n) = x) = p_X(x)$ per ogni $n$, quindi
\[ E[X(n)] = \sum_{x \in \mathcal{X}} x\, P(X(n) = x) = \sum_{x \in \mathcal{X}} x\, p_X(x) = \mu_X , \]
che non contiene $n$.
\passo{b) Dipende solo dalla differenza} Per ogni $h$, la stazionarietà del secondo ordine dice che le due pmf
congiunte sono uguali, e quindi anche le medie calcolate con esse:
\[ \begin{aligned}
   E[X(n_1) X(n_2)] &= \sum_{x_1 \in \mathcal{X}} \sum_{x_2 \in \mathcal{X}} x_1 x_2\, P(X(n_1) = x_1, X(n_2) = x_2) \\
   &= \sum_{x_1 \in \mathcal{X}} \sum_{x_2 \in \mathcal{X}} x_1 x_2\, P(X(n_1 + h) = x_1, X(n_2 + h) = x_2) =
   E[X(n_1 + h)\, X(n_2 + h)] .
   \end{aligned} \]
Scegliendo $h = -n_1$: $E[X(n_1) X(n_2)] = E[X(0)\, X(n_2 - n_1)]$, che è una funzione della sola differenza
$n_2 - n_1$.
\passo{b) Anzi, del valore assoluto} Per la proprietà commutativa del prodotto, $E[X(n_1) X(n_2)] = E[X(n_2) X(n_1)]$:
la stessa funzione dà lo stesso valore in $n_2 - n_1$ e in $n_1 - n_2$. Dipende quindi solo da $\abs{n_2 - n_1}$.
\end{dimostrazione}

Queste due proprietà sono così utili che si prendono come definizione di una stazionarietà più debole, che riguarda
solo medie e correlazioni.

\dfn{Autocorrelazione}{
  L'\term{autocorrelazione} di un processo $X(n)$ è la correlazione fra i suoi campioni in due istanti:
  \[ R_X(n_1, n_2) = E[X(n_1)\, X(n_2)] . \]
}

\dfn{Processo stazionario in senso lato}{
  Un processo, continuo o discreto, non necessariamente stazionario al secondo ordine, si dice \term{stazionario in
  senso lato} (\term{SSL}; in inglese \emph{wide-sense stationary}, WSS) se:
  \begin{itemize}
    \item la sua media non dipende dal tempo: $E[X(n)] = \mu_X$;
    \item la sua autocorrelazione dipende solo dalla differenza fra gli istanti:
          $R_X(n_1, n_2) = E[X(n_1) X(n_2)] = R_X(n_2 - n_1)$.
  \end{itemize}
}

Per il Teorema~\ref{teo:medie2} ogni processo stazionario del secondo ordine è SSL. Si scrive spesso
$R_X(k) = E[X(n) X(n + k)]$, con $k$ la distanza fra i campioni: per la commutatività del prodotto è una funzione pari,
$R_X(-k) = R_X(k)$, e $R_X(0) = E[X^2(n)]$ è il valore quadratico medio, cioè la \textbf{potenza media} del processo.

\warning{SSL non vuol dire stazionario}{
  La stazionarietà in senso lato riguarda solo medie e correlazioni, non le densità. Si prenda un processo a campioni
  indipendenti con $X(n) \sim \mathcal{N}(0, 1)$ negli istanti pari e $X(n) \sim \mathcal{U}(-\sqrt{3}, \sqrt{3})$ in
  quelli dispari. La media è sempre 0, la varianza è sempre 1 (per l'uniforme su $(-a, a)$ vale $\frac{(2a)^2}{12} =
  \frac{a^2}{3} = 1$), e $R_X(n_1, n_2) = 0$ per $n_1 \neq n_2$: il processo è SSL. Ma la densità del primo ordine
  cambia da un istante all'altro (una campana, poi un rettangolo): il processo non è stazionario nemmeno del primo
  ordine.
}

\subsection{Il vettore dei campioni di un processo SSL}

Cosa succede se si prelevano $M$ campioni da un processo SSL e li si mette in un vettore? La struttura della
stazionarietà si riflette in quella della matrice di covarianza.

\teorema[vettoressl]{Il vettore dei campioni di un processo SSL}{
  Sia $X(n)$ un processo SSL e sia $\vett{X} = [X_1, X_2, \ldots, X_M]^T$ il vettore dei suoi campioni negli istanti
  $n_1, n_2, \ldots, n_M$, cioè $X_i = X(n_i)$. Allora:
  \begin{enumerate}[label=\alph*)]
    \item il vettore delle medie ha tutti gli elementi uguali: $\vett{\mu}_{\vett{X}} = [\mu, \mu, \ldots, \mu]^T =
          \mu\, \vett{1}$, dove $\mu$ è uno scalare e $\vett{1} = [1, 1, \ldots, 1]^T$ è il vettore $M \times 1$ di tutti
          «1»;
    \item $E[X_i X_j] = f(\abs{n_i - n_j})$ per un'opportuna funzione $f$; in particolare il valore quadratico medio
          $E[X_i^2] = f(0)$ e la varianza $\sigma_X^2 = f(0) - \mu^2$ non dipendono da $i$;
    \item il coefficiente di correlazione dipende solo dalla distanza fra gli istanti, $\rho_{i,j} =
          \rho_{\abs{n_i - n_j}}$, e la matrice di covarianza è simmetrica, della forma
          \[ \matr{C}_{\vett{X}} = \sigma_X^2 \begin{bmatrix} 1 & \rho_{1,2} & \cdots & \rho_{1,M} \\ \rho_{2,1} & 1 &
             \cdots & \rho_{2,M} \\ \vdots & \vdots & \ddots & \vdots \\ \rho_{M,1} & \rho_{M,2} & \cdots & 1
             \end{bmatrix} ; \]
    \item se il passo di campionamento è costante ($n_{i+1} - n_i$ uguale per ogni $i$), $\matr{C}_{\vett{X}}$ è una
          \term{matrice di Toeplitz}: gli elementi di ciascuna diagonale sono tutti uguali.
  \end{enumerate}
}

\begin{dimostrazione}
\passo{a)} Per la prima condizione SSL, $E[X_i] = E[X(n_i)] = \mu$ per ogni $i$.
\passo{b)} Per la seconda condizione, $E[X_i X_j] = R_X(n_j - n_i)$. Siccome $E[X_i X_j] = E[X_j X_i] = R_X(n_i -
n_j)$, la funzione dà lo stesso valore in $n_j - n_i$ e nel suo opposto: dipende solo da $\abs{n_i - n_j}$. Con
$i = j$ la distanza è zero, quindi $E[X_i^2] = f(0)$, e $\sigma_{X_i}^2 = E[X_i^2] - E^2[X_i] = f(0) - \mu^2$.
\passo{c)} Tutte le deviazioni standard sono uguali a $\sigma_X$, e $COV(X_i, X_j) = E[X_i X_j] - \mu^2$ (Capitolo~9):
\[ \rho_{i,j} = \frac{COV(X_i, X_j)}{\sigma_{X_i}\, \sigma_{X_j}} = \frac{E[X_i X_j] - \mu^2}{\sigma_X^2} =
   \frac{E[X(n_i)\, X(n_j)] - \mu^2}{\sigma_X^2} = \rho_{\abs{n_i - n_j}} . \]
Ogni elemento della matrice di covarianza è $C_{X_i X_j} = \rho_{i,j}\, \sigma_X^2$; sulla diagonale $\rho_{i,i} = 1$,
e $\rho_{i,j} = \rho_{j,i}$ perché la distanza fra $n_i$ e $n_j$ è la stessa nei due versi.
\passo{d)} Con passo costante $\Delta$, $n_j - n_i = (j - i)\Delta$: l'elemento di posto $(i, j)$ dipende solo da
$\abs{j - i}$, cioè dalla diagonale su cui si trova.
\end{dimostrazione}

\begin{figure}[H]
\centering
\begin{tikzpicture}[ampersand replacement=\&, every left delimiter/.style={xshift=0.9ex},
  every right delimiter/.style={xshift=-0.9ex}]
  \matrix (T) [matrix of math nodes, left delimiter={[}, right delimiter={]},
    nodes={minimum width=2.6em, minimum height=2em, font=\small}] {
    |[fill=mspiA!30]| 1 \& |[fill=mspiB!25]| \rho_1 \& |[fill=mspiC!25]| \rho_2 \& |[fill=evid]| \rho_3 \&
      |[fill=mspiric!20]| \rho_4 \\
    |[fill=mspiB!25]| \rho_1 \& |[fill=mspiA!30]| 1 \& |[fill=mspiB!25]| \rho_1 \& |[fill=mspiC!25]| \rho_2 \&
      |[fill=evid]| \rho_3 \\
    |[fill=mspiC!25]| \rho_2 \& |[fill=mspiB!25]| \rho_1 \& |[fill=mspiA!30]| 1 \& |[fill=mspiB!25]| \rho_1 \&
      |[fill=mspiC!25]| \rho_2 \\
    |[fill=evid]| \rho_3 \& |[fill=mspiC!25]| \rho_2 \& |[fill=mspiB!25]| \rho_1 \& |[fill=mspiA!30]| 1 \&
      |[fill=mspiB!25]| \rho_1 \\
    |[fill=mspiric!20]| \rho_4 \& |[fill=evid]| \rho_3 \& |[fill=mspiC!25]| \rho_2 \& |[fill=mspiB!25]| \rho_1 \&
      |[fill=mspiA!30]| 1 \\
  };
  \node[left=0.4cm of T] {$\matr{C}_{\vett{X}} = \sigma_X^2$};
  \node[font=\small, anchor=west, text width=7.6cm] at ([xshift=0.8cm]T.east)
    {Cinque campioni consecutivi di un processo SSL. L'elemento $(i, j)$ dipende solo dalla distanza $\abs{i - j}$: la
     diagonale principale (distanza 0) contiene tutti 1, le diagonali vicine (distanza 1) tutti $\rho_1$, e così via.
     La matrice è simmetrica e \textbf{di Toeplitz}: bastano i numeri della prima riga, $1, \rho_1, \ldots, \rho_4$, per
     scriverla tutta.};
\end{tikzpicture}
\caption{La matrice di covarianza di campioni equispaziati di un processo SSL: ogni colore è una diagonale.}
\end{figure}

Si osservi che la matrice di covarianza $\matr{C}_{\vett{X}}$ è simmetrica anche se il vettore viene estratto da un
processo non SSL, perché la media del prodotto di due variabili aleatorie non dipende dall'ordine dei fattori. La
struttura, però, non sarà quella precedente: per esempio sulla diagonale principale potrebbero esserci valori diversi,
perché la varianza può dipendere dal tempo.

% ============================================================
\section{Quattro processi a confronto}
% ============================================================

Per rendere concrete le definizioni, studiamo quattro processi: tre stazionari in senso lato, con memorie diverse, e
uno che non lo è.

\begin{esempio}{Il processo di Bernoulli: nessuna memoria}
Per il processo di Bernoulli a valori $\pm 1$: $\mu = \frac12 (-1) + \frac12 (+1) = 0$, e $X^2(n) = 1$ sempre, quindi
$R_X(0) = 1$. Per $k \neq 0$ i due campioni sono indipendenti, e $R_X(k) = E[X(n)]\, E[X(n + k)] = 0$. Quindi
\[ R_X(k) = \begin{cases} 1 & k = 0 \\ 0 & k \neq 0 \end{cases} \qquad\text{e}\qquad \matr{C}_{\vett{X}} = \matr{I} \]
per qualsiasi vettore di campioni: un campione non dice nulla sugli altri. Il processo è SSL (anzi stazionario in senso
stretto, per il Teorema~\ref{teo:indipstaz}).
\end{esempio}

\begin{esempio}{La media mobile: un processo SSL con memoria}
Sia $X(n)$ il processo di Bernoulli a valori $\pm 1$ e sia $Y(n) = \frac12 [X(n) + X(n - 1)]$ la media di due campioni
consecutivi. $Y(n)$ vale $-1$ se entrambi i campioni valgono $-1$, $+1$ se valgono entrambi $+1$, e $0$ se sono diversi:
con probabilità $\frac14$, $\frac14$ e $\frac12$.
\begin{itemize}
  \item \textbf{Media}: $E[Y(n)] = \frac12 \left( E[X(n)] + E[X(n - 1)] \right) = 0$.
  \item \textbf{Autocorrelazione}: si sviluppa il prodotto e si usa $E[X(i)\, X(j)] = 1$ se $i = j$ e $0$ se $i \neq j$
        (esempio precedente):
        \[ R_Y(n, n + k) = \frac14 E\Big[ \left( X(n) + X(n - 1) \right) \left( X(n + k) + X(n + k - 1) \right) \Big]
           = \begin{cases} \frac14 (1 + 1) = \frac12 & k = 0 , \\ \frac14 & k = \pm 1 , \\ 0 & \abs{k} \geq 2 .
           \end{cases} \]
        Con $k = 0$ sopravvivono i due quadrati; con $k = 1$ solo il prodotto $X(n) \cdot X(n)$, perché $X(n)$ è
        l'unico campione in comune fra $Y(n)$ e $Y(n + 1)$ (con $k = -1$ è $X(n - 1)$); con $\abs{k} \geq 2$ le due
        medie non hanno campioni in comune.
\end{itemize}
$R_Y$ dipende solo da $k$: il processo è SSL, con $\sigma_Y^2 = \frac12$, $\rho_1 = \frac{1/4}{1/2} = \frac12$ e
$\rho_k = 0$ per $\abs{k} \geq 2$. Per quattro campioni consecutivi la matrice di covarianza è di Toeplitz:
\[ \matr{C}_{\vett{Y}} = \frac12 \begin{bmatrix} 1 & \frac12 & 0 & 0 \\ \frac12 & 1 & \frac12 & 0 \\ 0 & \frac12 & 1 &
   \frac12 \\ 0 & 0 & \frac12 & 1 \end{bmatrix} . \]
A differenza di $X(n)$, i campioni vicini di $Y(n)$ sono correlati: la media mobile «ricorda» il campione precedente e
rende la sequenza più lenta e meno frastagliata.
\begin{center}
\begin{tikzpicture}
\begin{groupplot}[group style={group size=1 by 2, vertical sep=0.5cm, xticklabels at=edge bottom}, realizzazione,
  width=15.5cm, height=3.2cm, xmin=-0.8, xmax=29.8, ymin=-1.4, ymax=1.4, ytick={-1,1}, xtick={0,5,10,15,20,25}]
\nextgroupplot
  \addplot[mspiA] coordinates {(0,-1) (1,-1) (2,-1) (3,1) (4,1) (5,1) (6,1) (7,1) (8,-1) (9,-1) (10,-1) (11,-1) (12,-1) (13,-1) (14,1) (15,-1) (16,-1) (17,1) (18,1) (19,-1) (20,1) (21,-1) (22,1) (23,1) (24,-1) (25,1) (26,1) (27,-1) (28,1) (29,-1)};
  \node[anchor=west, font=\small] at (axis cs:30.3,0) {$X(n)$};
\nextgroupplot[xlabel=$n$]
  \addplot[mspiB] coordinates {(0,0.0) (1,-1.0) (2,-1.0) (3,0.0) (4,1.0) (5,1.0) (6,1.0) (7,1.0) (8,0.0) (9,-1.0) (10,-1.0) (11,-1.0) (12,-1.0) (13,-1.0) (14,0.0) (15,0.0) (16,-1.0) (17,0.0) (18,1.0) (19,0.0) (20,0.0) (21,0.0) (22,0.0) (23,1.0) (24,0.0) (25,0.0) (26,1.0) (27,0.0) (28,0.0) (29,0.0)};
  \node[anchor=west, font=\small] at (axis cs:30.3,0) {$Y(n)$};
\end{groupplot}
\end{tikzpicture}
\end{center}
In alto la stessa realizzazione del processo di Bernoulli dell'esempio precedente; in basso la sua media mobile.
\end{esempio}

\begin{richiamo}{il prodotto di due coseni}
Dalle formule di addizione della trigonometria (che si studiano alle superiori)
\[ \cos(a - b) = \cos a \cos b + \sin a \sin b , \qquad \cos(a + b) = \cos a \cos b - \sin a \sin b . \]
Sommandole membro a membro, i termini con i seni si cancellano: $\cos(a - b) + \cos(a + b) = 2\cos a \cos b$, cioè
\[ \cos a \cos b = \frac12 \cos(a - b) + \frac12 \cos(a + b) . \]
Il prodotto di due coseni è la media di due coseni: uno con la \emph{differenza} degli argomenti e uno con la
\emph{somma}. Serve anche un secondo fatto: la media di un coseno su un numero intero di periodi è zero, perché le parti
sopra e sotto l'asse hanno la stessa area. Per esempio
\[ \int_{-1/2}^{1/2} \cos(4\pi u + c) \dd u = \left[ \frac{\sin(4\pi u + c)}{4\pi} \right]_{-1/2}^{1/2} =
   \frac{\sin(2\pi + c) - \sin(-2\pi + c)}{4\pi} = 0 , \]
perché il seno ha periodo $2\pi$ (con la regola della catena, la derivata di $\sin(4\pi u + c)$ è $4\pi \cos(4\pi u +
c)$, da cui la primitiva).
\end{richiamo}

\begin{esempio}{La sinusoide con fase a caso è SSL}
Riprendiamo $X(n) = A\cos(2\pi U + \varphi_n)$, con $U \sim \mathcal{U}\left( -\frac12, \frac12 \right)$ e
$\varphi_n = 2\pi f_0 n$.
\begin{itemize}
  \item \textbf{Media}: $E[X(n)] = \int_{-1/2}^{1/2} A\cos(2\pi u + \varphi_n) \dd u = 0$ (un periodo intero di
        coseno), qualunque sia $n$.
  \item \textbf{Autocorrelazione}: con la formula del richiamo,
        \[ X(n_1)\, X(n_2) = A^2 \cos(2\pi U + \varphi_{n_1}) \cos(2\pi U + \varphi_{n_2}) = \frac{A^2}{2}
           \cos(\varphi_{n_2} - \varphi_{n_1}) + \frac{A^2}{2} \cos(4\pi U + \varphi_{n_1} + \varphi_{n_2}) . \]
        Il primo termine non contiene $U$: è una costante. Il secondo, mediato su $U$, dà zero (due periodi interi di
        coseno, come nel richiamo). Quindi
        \[ R_X(n_1, n_2) = \frac{A^2}{2} \cos(\varphi_{n_2} - \varphi_{n_1}) = \frac{A^2}{2} \cos\left( 2\pi f_0 (n_2 -
           n_1) \right) , \]
        che dipende solo da $n_2 - n_1$ (e, essendo il coseno pari, solo da $\abs{n_2 - n_1}$).
\end{itemize}
Il processo è SSL, con $\mu = 0$ e $R_X(k) = \frac{A^2}{2} \cos(2\pi f_0 k)$. La potenza media $R_X(0) = \frac{A^2}{2}$
è il valore quadratico medio di una sinusoide di ampiezza $A$; e due campioni distanti un periodo ($k = 1/f_0$) hanno di
nuovo correlazione massima, perché la sinusoide si ripete. A differenza della media mobile, la memoria di questo processo
non si esaurisce mai: conoscendo due campioni si ricostruisce tutta la sinusoide.
\end{esempio}

\begin{figure}[H]
\centering
\begin{tikzpicture}
\begin{groupplot}[group style={group size=3 by 1, horizontal sep=1cm}, realizzazione, width=6.6cm, height=4.6cm,
  xmin=-8.6, xmax=8.6, ymin=-0.65, ymax=1.15, xtick={-8,-4,0,4,8}, ytick={-0.5,0.5,1},
  yticklabels={$-\frac12$,$\frac12$,$1$}, xlabel=$k$, every axis plot/.append style={mark size=1.4pt, line width=1pt}]
\nextgroupplot[title={Bernoulli}]
  \addplot[mspiA] coordinates {(-8,0) (-7,0) (-6,0) (-5,0) (-4,0) (-3,0) (-2,0) (-1,0) (0,1) (1,0) (2,0) (3,0) (4,0)
    (5,0) (6,0) (7,0) (8,0)};
\nextgroupplot[title={media mobile}]
  \addplot[mspiB] coordinates {(-8,0) (-7,0) (-6,0) (-5,0) (-4,0) (-3,0) (-2,0) (-1,0.25) (0,0.5) (1,0.25) (2,0) (3,0)
    (4,0) (5,0) (6,0) (7,0) (8,0)};
\nextgroupplot[title={sinusoide con fase a caso}]
  \addplot[mspiC, samples at={-8,-7,...,8}] {0.5*cos(deg(2*pi*x/12))};
  \addplot[mspiC!60, line width=0.5pt, dashed, no marks, domain=-8.4:8.4, samples=100, sharp plot]
    {0.5*cos(deg(2*pi*x/12))};
\end{groupplot}
\end{tikzpicture}
\caption{Le autocorrelazioni $R(k)$ dei tre processi SSL degli esempi (con $A = 1$ e $f_0 = \frac{1}{12}$ per la
sinusoide). Sono tutte pari e hanno il massimo in $k = 0$, la potenza media. La memoria del processo si legge nella
larghezza: nulla per Bernoulli, un passo per la media mobile, infinita per la sinusoide.}
\end{figure}

\begin{esempio}{La passeggiata aleatoria non è stazionaria}
Una pedina parte dalla casella 0 e a ogni istante si sposta di una casella avanti o indietro, a seconda dell'esito di
una moneta: la sua posizione dopo $n$ passi è $S(n) = X(1) + X(2) + \cdots + X(n)$, con $S(0) = 0$, dove $X(n)$ è il
processo di Bernoulli a valori $\pm 1$. È la \term{passeggiata aleatoria}.
\begin{itemize}
  \item \textbf{Media}: $E[S(n)] = \sum_{i=1}^{n} E[X(i)] = 0$, costante.
  \item \textbf{Varianza}: i passi sono indipendenti, quindi le varianze si sommano (Capitolo~9):
        $\sigma_{S(n)}^2 = \sum_{i=1}^{n} \sigma_{X(i)}^2 = n$. Cresce con il tempo.
  \item \textbf{Autocorrelazione}: per $n_1 \leq n_2$, $E[S(n_1)\, S(n_2)] = \sum_{i=1}^{n_1} \sum_{j=1}^{n_2} E[X(i)\,
        X(j)]$, e sopravvivono solo i termini con $i = j$, che sono $n_1$: $R_S(n_1, n_2) = \min(n_1, n_2)$.
\end{itemize}
La media non dipende dal tempo, ma l'autocorrelazione non dipende solo dalla differenza $n_2 - n_1$: per esempio
$R_S(1, 2) = 1$ e $R_S(10, 11) = 10$. Il processo \textbf{non} è SSL, e quindi nemmeno stazionario del secondo ordine;
anzi, non lo è neppure del primo ordine, perché la pmf di $S(n)$ si allarga con $n$. Nella figura le realizzazioni si
allontanano dall'origine come $\sqrt{n}$, la deviazione standard dopo $n$ passi.
\begin{center}
\begin{tikzpicture}
\begin{axis}[width=15.5cm, height=7cm, xmin=0, xmax=100, ymin=-25, ymax=25, xtick={0,20,40,60,80,100},
  ytick={-20,-10,0,10,20}, xlabel=$n$, ylabel={$S(n)$}, axis lines=left, tick label style={font=\scriptsize},
  label style={font=\small}, grid=major, grid style={mspigray!20}, legend style={font=\scriptsize, at={(0.5,1.03)},
  anchor=south, legend columns=2, draw=none}]
  \addplot[black, dashed, line width=0.8pt, domain=0:100, samples=101] {sqrt(x)};
  \addplot[black, dotted, line width=1pt, domain=0:100, samples=101] {2*sqrt(x)};
  \addplot[black, dashed, line width=0.8pt, domain=0:100, samples=101, forget plot] {-sqrt(x)};
  \addplot[black, dotted, line width=1pt, domain=0:100, samples=101, forget plot] {-2*sqrt(x)};
  \addplot[mspiA, line width=0.7pt, forget plot] coordinates {(0,0) (1,-1) (2,-2) (3,-3) (4,-2) (5,-3) (6,-2) (7,-1) (8,-2) (9,-3) (10,-4) (11,-3) (12,-2) (13,-1) (14,0) (15,1) (16,0) (17,-1) (18,0) (19,1) (20,2) (21,3) (22,4) (23,3) (24,2) (25,1) (26,0) (27,-1) (28,-2) (29,-1) (30,-2) (31,-3) (32,-2) (33,-3) (34,-2) (35,-1) (36,0) (37,1) (38,2) (39,3) (40,2) (41,3) (42,4) (43,3) (44,4) (45,5) (46,6) (47,7) (48,8) (49,9) (50,10) (51,11) (52,12) (53,11) (54,12) (55,11) (56,12) (57,13) (58,14) (59,15) (60,16) (61,17) (62,16) (63,17) (64,18) (65,17) (66,18) (67,17) (68,16) (69,15) (70,14) (71,13) (72,14) (73,13) (74,12) (75,11) (76,12) (77,11) (78,10) (79,9) (80,10) (81,9) (82,8) (83,9) (84,10) (85,9) (86,8) (87,7) (88,6) (89,5) (90,4) (91,3) (92,2) (93,1) (94,2) (95,3) (96,2) (97,1) (98,0) (99,-1) (100,0)};
  \addplot[mspiB, line width=0.7pt, forget plot] coordinates {(0,0) (1,-1) (2,-2) (3,-3) (4,-4) (5,-5) (6,-4) (7,-5) (8,-4) (9,-3) (10,-2) (11,-3) (12,-2) (13,-1) (14,-2) (15,-1) (16,0) (17,-1) (18,0) (19,-1) (20,-2) (21,-1) (22,-2) (23,-3) (24,-2) (25,-3) (26,-2) (27,-1) (28,0) (29,-1) (30,0) (31,1) (32,2) (33,3) (34,2) (35,1) (36,0) (37,1) (38,0) (39,-1) (40,-2) (41,-3) (42,-4) (43,-3) (44,-4) (45,-5) (46,-6) (47,-7) (48,-6) (49,-7) (50,-6) (51,-7) (52,-8) (53,-7) (54,-6) (55,-5) (56,-4) (57,-5) (58,-6) (59,-7) (60,-6) (61,-5) (62,-6) (63,-7) (64,-8) (65,-9) (66,-10) (67,-11) (68,-12) (69,-13) (70,-14) (71,-15) (72,-16) (73,-17) (74,-16) (75,-15) (76,-14) (77,-13) (78,-14) (79,-15) (80,-14) (81,-13) (82,-14) (83,-15) (84,-14) (85,-15) (86,-16) (87,-15) (88,-14) (89,-15) (90,-16) (91,-15) (92,-14) (93,-13) (94,-12) (95,-13) (96,-12) (97,-13) (98,-14) (99,-13) (100,-14)};
  \addplot[mspiC, line width=0.7pt, forget plot] coordinates {(0,0) (1,-1) (2,0) (3,-1) (4,0) (5,1) (6,2) (7,1) (8,0) (9,-1) (10,0) (11,-1) (12,0) (13,1) (14,0) (15,-1) (16,0) (17,-1) (18,-2) (19,-1) (20,0) (21,1) (22,0) (23,1) (24,2) (25,1) (26,2) (27,1) (28,0) (29,-1) (30,0) (31,1) (32,0) (33,-1) (34,-2) (35,-1) (36,-2) (37,-1) (38,0) (39,1) (40,2) (41,1) (42,2) (43,3) (44,4) (45,5) (46,6) (47,7) (48,8) (49,9) (50,8) (51,7) (52,8) (53,9) (54,10) (55,9) (56,8) (57,7) (58,8) (59,7) (60,6) (61,7) (62,6) (63,5) (64,4) (65,3) (66,4) (67,5) (68,4) (69,5) (70,6) (71,5) (72,6) (73,5) (74,4) (75,5) (76,4) (77,5) (78,4) (79,5) (80,4) (81,3) (82,4) (83,3) (84,4) (85,5) (86,4) (87,5) (88,6) (89,5) (90,6) (91,5) (92,6) (93,5) (94,6) (95,7) (96,8) (97,9) (98,10) (99,11) (100,12)};
  \addplot[mspiric, line width=0.7pt, forget plot] coordinates {(0,0) (1,1) (2,0) (3,-1) (4,-2) (5,-1) (6,0) (7,-1) (8,0) (9,1) (10,2) (11,1) (12,0) (13,-1) (14,0) (15,-1) (16,0) (17,-1) (18,-2) (19,-1) (20,-2) (21,-1) (22,-2) (23,-1) (24,-2) (25,-3) (26,-2) (27,-1) (28,0) (29,-1) (30,0) (31,-1) (32,-2) (33,-1) (34,-2) (35,-1) (36,-2) (37,-1) (38,0) (39,1) (40,2) (41,3) (42,2) (43,1) (44,2) (45,1) (46,2) (47,3) (48,2) (49,1) (50,0) (51,1) (52,2) (53,3) (54,4) (55,3) (56,2) (57,3) (58,4) (59,3) (60,4) (61,3) (62,4) (63,5) (64,4) (65,5) (66,6) (67,5) (68,6) (69,7) (70,6) (71,7) (72,8) (73,9) (74,10) (75,11) (76,12) (77,11) (78,10) (79,11) (80,10) (81,9) (82,10) (83,9) (84,8) (85,9) (86,8) (87,9) (88,8) (89,7) (90,6) (91,7) (92,8) (93,7) (94,8) (95,9) (96,8) (97,9) (98,10) (99,11) (100,10)};
  \addplot[mspiA!50, line width=0.7pt, forget plot] coordinates {(0,0) (1,1) (2,2) (3,3) (4,4) (5,3) (6,4) (7,3) (8,4) (9,5) (10,4) (11,3) (12,2) (13,3) (14,2) (15,3) (16,2) (17,3) (18,2) (19,1) (20,0) (21,1) (22,2) (23,3) (24,4) (25,5) (26,6) (27,7) (28,8) (29,7) (30,8) (31,9) (32,10) (33,11) (34,10) (35,11) (36,12) (37,13) (38,14) (39,13) (40,12) (41,13) (42,14) (43,13) (44,12) (45,11) (46,12) (47,11) (48,12) (49,11) (50,12) (51,11) (52,12) (53,11) (54,12) (55,13) (56,14) (57,13) (58,14) (59,15) (60,14) (61,15) (62,14) (63,13) (64,12) (65,13) (66,14) (67,13) (68,14) (69,15) (70,16) (71,15) (72,16) (73,15) (74,14) (75,13) (76,12) (77,13) (78,14) (79,13) (80,12) (81,11) (82,10) (83,11) (84,12) (85,13) (86,12) (87,13) (88,12) (89,13) (90,14) (91,13) (92,12) (93,13) (94,14) (95,15) (96,16) (97,15) (98,16) (99,17) (100,18)};
  \addplot[mspiB!55, line width=0.7pt, forget plot] coordinates {(0,0) (1,-1) (2,0) (3,1) (4,2) (5,1) (6,0) (7,1) (8,0) (9,-1) (10,-2) (11,-3) (12,-2) (13,-1) (14,0) (15,-1) (16,0) (17,-1) (18,0) (19,-1) (20,-2) (21,-3) (22,-4) (23,-3) (24,-4) (25,-3) (26,-2) (27,-3) (28,-2) (29,-3) (30,-4) (31,-3) (32,-4) (33,-5) (34,-6) (35,-7) (36,-8) (37,-9) (38,-10) (39,-9) (40,-10) (41,-11) (42,-10) (43,-9) (44,-10) (45,-11) (46,-12) (47,-13) (48,-14) (49,-13) (50,-14) (51,-15) (52,-16) (53,-15) (54,-14) (55,-13) (56,-14) (57,-15) (58,-16) (59,-17) (60,-18) (61,-17) (62,-18) (63,-17) (64,-16) (65,-15) (66,-16) (67,-15) (68,-16) (69,-15) (70,-14) (71,-13) (72,-14) (73,-15) (74,-14) (75,-13) (76,-12) (77,-11) (78,-10) (79,-11) (80,-10) (81,-11) (82,-12) (83,-13) (84,-14) (85,-15) (86,-16) (87,-15) (88,-16) (89,-17) (90,-18) (91,-17) (92,-16) (93,-17) (94,-18) (95,-19) (96,-18) (97,-19) (98,-20) (99,-21) (100,-22)};
  \addplot[mspiC!55, line width=0.7pt, forget plot] coordinates {(0,0) (1,-1) (2,-2) (3,-1) (4,0) (5,-1) (6,-2) (7,-1) (8,0) (9,-1) (10,0) (11,1) (12,0) (13,1) (14,2) (15,1) (16,2) (17,1) (18,0) (19,1) (20,2) (21,3) (22,4) (23,3) (24,4) (25,3) (26,2) (27,3) (28,4) (29,5) (30,6) (31,7) (32,8) (33,7) (34,8) (35,9) (36,8) (37,9) (38,10) (39,11) (40,10) (41,9) (42,10) (43,11) (44,10) (45,9) (46,10) (47,9) (48,10) (49,9) (50,8) (51,7) (52,6) (53,7) (54,6) (55,7) (56,8) (57,9) (58,10) (59,9) (60,10) (61,11) (62,10) (63,11) (64,10) (65,9) (66,8) (67,7) (68,8) (69,7) (70,8) (71,9) (72,8) (73,9) (74,10) (75,9) (76,8) (77,9) (78,8) (79,7) (80,8) (81,9) (82,10) (83,9) (84,8) (85,7) (86,8) (87,7) (88,6) (89,7) (90,8) (91,7) (92,8) (93,9) (94,10) (95,11) (96,12) (97,13) (98,12) (99,11) (100,10)};
  \addplot[mspigray, line width=0.7pt, forget plot] coordinates {(0,0) (1,-1) (2,-2) (3,-3) (4,-4) (5,-3) (6,-2) (7,-1) (8,-2) (9,-1) (10,-2) (11,-1) (12,0) (13,1) (14,0) (15,1) (16,0) (17,1) (18,0) (19,-1) (20,0) (21,-1) (22,-2) (23,-3) (24,-4) (25,-3) (26,-4) (27,-5) (28,-6) (29,-5) (30,-6) (31,-7) (32,-8) (33,-7) (34,-8) (35,-9) (36,-8) (37,-7) (38,-6) (39,-7) (40,-6) (41,-5) (42,-6) (43,-5) (44,-6) (45,-5) (46,-4) (47,-3) (48,-2) (49,-1) (50,-2) (51,-3) (52,-4) (53,-3) (54,-2) (55,-3) (56,-2) (57,-3) (58,-4) (59,-3) (60,-2) (61,-1) (62,-2) (63,-3) (64,-4) (65,-5) (66,-6) (67,-5) (68,-6) (69,-5) (70,-6) (71,-7) (72,-8) (73,-9) (74,-8) (75,-9) (76,-8) (77,-9) (78,-10) (79,-11) (80,-12) (81,-13) (82,-14) (83,-15) (84,-14) (85,-15) (86,-14) (87,-15) (88,-16) (89,-15) (90,-14) (91,-15) (92,-16) (93,-15) (94,-16) (95,-15) (96,-14) (97,-15) (98,-16) (99,-15) (100,-14)};
  \legend{{$\pm\sqrt{n}$ (una deviazione standard)}, {$\pm 2\sqrt{n}$}}
\end{axis}
\end{tikzpicture}
\end{center}
Otto realizzazioni della passeggiata aleatoria (simulazione): partono tutte da 0 e si sparpagliano sempre di più.
\end{esempio}

\gbox{In sintesi}{mspiA!80!black}{
\begin{itemize}
  \item \textbf{Regola della catena}: $p_{\vett{X}}(\vett{x}) = \prod_{i} p_{X_i \mid X_1, \ldots, X_{i-1}}(x_i \mid
        x_1, \ldots, x_{i-1})$ (lo stesso con le densità): la congiunta si costruisce una variabile alla volta.
  \item Un \textbf{processo aleatorio a tempo discreto} associa a ogni $\omega$ una successione $\{X(n, \omega)\}_{n \in
        \Z}$: fissato $\omega$ si ha una realizzazione, fissato $n$ una variabile aleatoria.
  \item \textbf{Caratterizzazione}: al primo ordine $f_{X(n)}(x; n)$, al secondo $f_{X(n_1), X(n_2)}(x_1, x_2; n_1,
        n_2)$, completa con le densità di ogni vettore di $M$ campioni. \textbf{Stazionarietà} (di ordine $M$): le
        densità non cambiano traslando rigidamente gli istanti; l'ordine $M$ implica gli ordini più bassi.
  \item \textbf{Processo indipendente}: la densità di $M$ campioni è il prodotto delle marginali; se è anche stazionario
        basta una sola pmf (o PDF). Esempi: Bernoulli $\pm 1$, quaternario.
  \item \textbf{SSL}: media costante e autocorrelazione $R_X(n_1, n_2) = E[X(n_1)X(n_2)] = R_X(n_2 - n_1)$, pari in
        $k = n_2 - n_1$. Stazionario del secondo ordine $\Rightarrow$ SSL, non viceversa. I campioni equispaziati di un
        processo SSL hanno $\vett{\mu}_{\vett{X}} = \mu\vett{1}$ e una matrice di covarianza simmetrica e di Toeplitz,
        $\sigma_X^2 [\rho_{\abs{i - j}}]$.
\end{itemize}
}
